%%%%%%%%%%%%%%%%%
% CV tailored for ADEME - Chef.fe de projet Transport PRODICEE
% Positioning: Chef de projet Transport / Transition énergétique | PRODICEE, coordination, KPI, données & acteurs publics
%%%%%%%%%%%%%%%%%

\documentclass[8pt,a4paper,ragged2e,withhyper]{altacv}

\geometry{left=1.25cm,right=1.25cm,top=.5cm,bottom=1.cm,columnsep=1.2cm}

\usepackage{paracol}
\usepackage{fontawesome5}
\usepackage{hyperref}

\iftutex
  \setmainfont{Roboto Slab}
  \setsansfont{Lato}
  \renewcommand{\familydefault}{\sfdefault}
\else
  \usepackage[rm]{roboto}
  \usepackage[defaultsans]{lato}
  \renewcommand{\familydefault}{\sfdefault}
\fi

\definecolor{SlateGrey}{HTML}{2E2E2E}
\definecolor{LightGrey}{HTML}{666666}
\definecolor{DarkPastelRed}{HTML}{8AC0D9}
\definecolor{PastelRed}{HTML}{00B0DA}
\definecolor{GoldenEarth}{HTML}{B4AD0C}

\colorlet{name}{black}
\colorlet{tagline}{PastelRed}
\colorlet{heading}{DarkPastelRed}
\colorlet{headingrule}{GoldenEarth}
\colorlet{subheading}{PastelRed}
\colorlet{accent}{PastelRed}
\colorlet{emphasis}{SlateGrey}
\colorlet{body}{LightGrey}

\definecolor{violet}{HTML}{5A2582}
\definecolor{rouge}{HTML}{F53B3B}
\definecolor{noir}{HTML}{00232C}
\definecolor{bleu}{HTML}{00B0DA}
\definecolor{rose}{HTML}{F831A0}
\definecolor{jaune}{HTML}{FFEC00}
\definecolor{vert}{HTML}{34AD6C}

\renewcommand{\namefont}{\Huge\rmfamily\bfseries}
\renewcommand{\personalinfofont}{\footnotesize}
\renewcommand{\cvsectionfont}{\LARGE\rmfamily\bfseries}
\renewcommand{\cvsubsectionfont}{\large\bfseries}
\renewcommand{\cvItemMarker}{{\small\textbullet}}
\renewcommand{\cvRatingMarker}{\faCircle}

\input{../../_shared/pubs-num.tex}

% auto_tex_manager layout tuning
\emergencystretch=3em
\hfuzz=60pt
\hbadness=9999
\sloppy

\begin{document}

\name{BOILEAU Guillaume}
\tagline{Chef de projet Transport / Transition énergétique | Coordination, KPI, données, acteurs publics \& pilotage}

\photoL{2.8cm}{../../_shared/moi}

\personalinfo{%
  \email{guillaume.boileaupro@gmail.com}
  \phone{+33 6 59 94 85 84}
  \location{Alpes-Maritimes, FRANCE}
  \linkedin{guillaume-boileau-891b81265}
  \researchgate{Guillaume-Boileau-2}
  \orcid{0000-0002-3576-6968}
}

\makecvheader
\vspace{-0.3cm}

\AtBeginEnvironment{itemize}{\small}
\columnratio{0.64}

\begin{paracol}{2}

\cvsection{Experience}

\cvevent{\textbf{Software Engineer -- Mobilité électrique, données opérationnelles \& support terrain}}{VEV Platform Services France}{2025 -- 2026}{Valbonne, France}

\textbf{Contexte :} Plateforme CPMS liée à des infrastructures de recharge de véhicules électriques, avec services applicatifs, flux de données, contraintes terrain, exploitation opérationnelle et enjeux de fiabilisation.

\textbf{Objectif :} Contribuer à la fiabilisation de services numériques connectés à des infrastructures de mobilité électrique, en analysant les flux, les incidents, les comportements applicatifs et les besoins opérationnels.

\begin{itemize}
  \item \textbf{Mobilité électrique :} Travail dans l’écosystème des infrastructures de recharge, à l’interface entre logiciel, données opérationnelles et contraintes terrain.
  \item \textbf{Données \& exploitation :} Analyse de flux FTP, interruptions de flux, incohérences de données et comportements inattendus.
  \item \textbf{Support opérationnel :} Diagnostic d’incidents, analyse des causes, formalisation de constats et contribution aux actions de correction.
  \item \textbf{Validation terrain :} Vérification de la cohérence entre données applicatives, comportement des équipements et attendus fonctionnels.
  \item \textbf{Documentation technique :} Rédaction de constats, analyses d’anomalies et éléments de suivi utiles aux équipes.
  \item \textbf{Développement logiciel :} Contribution à des composants TypeScript, Angular et Node.js intégrés dans une plateforme opérationnelle.
  \item \textbf{Coordination :} Suivi collaboratif des sujets via Git, Linux, Zenhub et échanges avec les équipes techniques.
\end{itemize}

\divider

\cvevent{\textbf{Ingénieur R\&D senior -- Coordination, données, validation \& programme scientifique}}{CNES}{2023 -- 2025}{Nice, France}

\textbf{Contexte :} Activités R\&D et développement logiciel pour le segment sol de la mission spatiale LISA ESA/NASA, avec chaînes de données mission L0--L3, simulation, intégration de modèles, validation technique et environnement CNES/ESA.

\textbf{Objectif :} Structurer, suivre, documenter et valider des travaux complexes impliquant plusieurs contributeurs, données hétérogènes, méthodes de simulation, indicateurs de qualité et revues collaboratives.

\begin{itemize}
  \item \textbf{Pilotage de travaux complexes :} Structuration des besoins, hypothèses, dépendances, jalons, résultats attendus et points de vigilance.
  \item \textbf{Coordination transverse :} Interface avec contributeurs scientifiques, logiciels et institutionnels dans un environnement international.
  \item \textbf{Données \& indicateurs :} Structuration de données, métriques, produits d’analyse, contrôles de cohérence et éléments de reporting.
  \item \textbf{Validation de résultats :} Définition de critères d’acceptabilité, méthodes de comparaison, analyse d’écarts et synthèse des résultats.
  \item \textbf{Documentation \& revues :} Production de supports techniques, comptes rendus, synthèses, hypothèses, limites et recommandations.
  \item \textbf{Analyse de risques techniques :} Identification d’écarts, dépendances, incertitudes, impacts et actions correctrices.
  \item \textbf{Plateforme Python :} Développement de CATpyLISA pour l’intégration, la comparaison, la validation et la revue technique de modèles.
  \textcolor{DarkPastelRed}{\href{https://gitlab.oca.eu/gboileau/lisa_l3_tools}{\bf GitLab: CATpyLISA}}
\end{itemize}

\divider

\cvevent{\textbf{Data Scientist -- Données capteurs, performance \& recommandations techniques}}{Universiteit Antwerpen, Belgium}{2021 -- 2023}{Antwerp, Belgium}

\textbf{Contexte :} Travaux de data science appliqués à des instruments de précision, combinant données capteurs, bruit environnemental, modélisation statistique, machine learning, validation et recommandations techniques.

\textbf{Objectif :} Développer des workflows robustes pour analyser des données complexes, comparer des configurations, évaluer des performances et produire des recommandations exploitables.

\begin{itemize}
  \item \textbf{Analyse de données :} Exploitation de mesures capteurs, corrélations spatiales, propagation, bruit et impacts sur la performance.
  \item \textbf{KPI \& performance :} Figures de mérite, analyses de sensibilité, contrôles de cohérence et comparaison de configurations.
  \item \textbf{Machine learning :} Approches de réduction de bruits non stationnaires avec canaux témoins, machine learning et deep learning.
  \item \textbf{Reporting technique :} Production de synthèses, recommandations, visualisations et éléments de décision.
  \item \textbf{Coordination internationale :} Échanges avec parties prenantes scientifiques et techniques dans un contexte multiculturel.
  \item \textbf{Autonomie \& rigueur :} Gestion de sujets complexes, méthodes reproductibles, documentation et suivi des résultats.
\end{itemize}

\divider

\cvevent{\textbf{Ingénieur R\&D junior -- Simulation, signal \& rédaction scientifique}}{ARTEMIS, Observatoire de la Côte d’Azur}{2018 -- 2021}{Nice, France}

\textbf{Contexte :} Recherche appliquée liée à l’analyse de signaux faibles, à la simulation numérique, au traitement du signal et à la préparation de missions spatiales.

\textbf{Objectif :} Développer des méthodes d’analyse, automatiser des simulations et produire des résultats scientifiques documentés et exploitables.

\begin{itemize}
  \item \textbf{Simulation numérique :} Développement d’approches de simulation pour environnements complexes et grands jeux de données.
  \item \textbf{Analyse \& synthèse :} Traitement de signaux faibles, comparaison de résultats, études de sensibilité et validation.
  \item \textbf{Rédaction :} Contribution à des publications, rapports, supports de présentation et documentation scientifique.
  \item \textbf{Restitution :} Présentation de résultats complexes à des publics scientifiques et techniques.
\end{itemize}

\switchcolumn

\cvsection{Profile}

\textbf{Docteur en physique et ingénieur R\&D / data avec expérience en pilotage de travaux complexes, coordination transverse, analyse de données, KPI, documentation, synthèse et environnement mobilité électrique.}

Positionnement pour ce rôle : contribuer au pilotage du volet Transport de PRODICEE en apportant une capacité forte d’analyse, de coordination, de structuration des données, de suivi d’actions, de reporting, de rédaction et d’animation technique.

\begin{itemize}
  \item Expérience concrète liée à la mobilité électrique et aux infrastructures de recharge chez VEV.
  \item Forte capacité d’analyse de données, indicateurs, flux, résultats techniques et anomalies.
  \item Expérience en coordination de contributeurs, suivi de livrables, documentation, synthèse et restitution.
  \item Culture institutionnelle et R\&D : CNES/ESA, revues techniques, méthodes, validation, traçabilité et travail avec parties prenantes.
  \item Montée ciblée sur le dispositif CEE, les acteurs nationaux transport, CEREMA/ATEE et l’évaluation technico-économique des fiches CEE.
\end{itemize}

\cvsection{Languages}

\cvskill{English}{4}
\divider
\cvskill{Spanish}{2}
\divider

\medskip

\cvsection{Education}

\cvevent{PhD in Physics}{Université Côte d’Azur}{2018 -- 2021}{Nice, France}

\divider

\cvevent{Selective Graduate Program in Fundamental Physics (Magistère)}{Université Paris-Saclay}{2016 -- 2018}{Orsay, France}

\divider

\cvevent{MSc in Astronomy and Astrophysics}{Observatoire de Paris / Université Paris-Saclay}{2017 -- 2018}{Paris, France}

\divider

\cvevent{BSc in Fundamental Physics}{Université Paris-Saclay}{2015 -- 2016}{Orsay, France}

\cvsection{Professional Training}

\cvevent{Machine Learning in Python with scikit-learn}{FUN MOOC -- Inria}{2026}{France Université Numérique}

\small{
Predictive modelling, supervised learning, model evaluation, cross-validation, metrics, pipelines and practical machine learning workflows with Python and scikit-learn.
}

\divider

\cvevent{Recherche reproductible : principes méthodologiques pour une science transparente}{FUN MOOC -- Inria}{2025}{France Université Numérique}

\small{
Reproducible research, GitLab, notebooks/Jupyter, data management, computational documentation, software environments and reproducibility methods.
}

\divider

\cvevent{Continuous Professional Training -- Software, Data, AI and Systems}{OpenClassrooms}{2024 -- 2026}{}

\small{
Python, SQL, Machine Learning, AI, Linux, JavaScript, TypeScript, Angular, Node.js, Django, REST APIs, data analysis, software engineering and algorithms.
}

\end{paracol}

\cvsection{Projects}

\cvevent{VEV -- Mobilité électrique, données opérationnelles \& infrastructures de recharge}{CPMS / EV charging / Data flows}{2025 -- 2026}{}

Contribution à une plateforme numérique connectée à des infrastructures de recharge de véhicules électriques, avec analyse de flux, incidents, données opérationnelles et contraintes terrain.

\textbf{Rôle :} analyse de flux, support technique, investigation d’incidents, validation de comportements applicatifs, documentation et contribution à la fiabilisation.

\textbf{Apport ADEME / Transport :} compréhension opérationnelle d’un environnement lié à la mobilité électrique, aux données d’usage, aux infrastructures et à la fiabilité des services numériques.

\divider

\cvevent{CATpyLISA -- Coordination R\&D, données et validation}{Mission spatiale LISA ESA/NASA -- CNES/ESA}{2023 -- 2025}{}

Plateforme Python dédiée à l’intégration de modèles, à la structuration de données, à la comparaison de résultats, à la validation technique et à la revue collaborative de workflows scientifiques.

\textbf{Rôle :} cadrage, développement Python, structuration de données, validation, analyse d’écarts, documentation technique, synthèse de résultats et coordination avec plusieurs contributeurs.

\textbf{Apport programme :} méthode de pilotage technique, suivi de résultats, analyse de risques, production de synthèses et coordination de travaux complexes.

\divider

\cvevent{Workflows data, KPI et recommandations techniques}{Capteurs / Performance / Modélisation}{2021 -- 2023}{}

Développement de workflows de données pour analyser des mesures complexes, produire des indicateurs, comparer des configurations et formuler des recommandations techniques.

\textbf{Rôle :} analyse de données, production d’indicateurs, validation de résultats, reporting, visualisation et restitution de conclusions.

\textbf{Apport Transport / CEE :} capacité à croiser des données, évaluer des impacts, produire des KPI et documenter des résultats pour éclairer la décision.

\cvsection{Compétences clés}

\vspace{-.3cm}

\cvsubsection{Pilotage programme \& coordination}

{\small
\cvtag{Pilotage de projet}
\cvtag{Coordination transverse}
\cvtag{Animation de réseau -- ramp-up}
\cvtag{Groupes de travail}
\cvtag{Comités techniques}
\cvtag{COMEX / COTECH -- awareness}
\cvtag{Planning}
\cvtag{Suivi d’actions}
\cvtag{Pilotage prestataires -- ramp-up}
\cvtag{Travail en autonomie}
\cvtag{Management de projet}
\cvtag{Diplomatie}
}

\divider\smallskip
\vspace{-.3cm}

\cvsubsection{Transport, mobilité \& transition énergétique}

{\small
\cvtag{Mobilité électrique}
\cvtag{Infrastructures de recharge}
\cvtag{Transport -- ramp-up}
\cvtag{Transition énergétique}
\cvtag{CEE -- ramp-up}
\cvtag{Fiches CEE -- ramp-up}
\cvtag{Économies d’énergie -- awareness}
\cvtag{Incitativité -- ramp-up}
\cvtag{CEREMA -- awareness}
\cvtag{ATEE -- awareness}
\cvtag{Politiques publiques -- awareness}
}

\divider\smallskip
\vspace{-.3cm}

\cvsubsection{Données, KPI \& évaluation}

{\small
\cvtag{Analyse de données}
\cvtag{KPI}
\cvtag{Suivi technique}
\cvtag{Suivi financier -- ramp-up}
\cvtag{Budget -- ramp-up}
\cvtag{Enquêtes terrain -- ramp-up}
\cvtag{Croisement de données}
\cvtag{Reporting}
\cvtag{Tableaux de bord}
\cvtag{Power BI}
\cvtag{Python}
\cvtag{SQL}
\cvtag{pandas}
}

\divider\smallskip
\vspace{-.3cm}

\cvsubsection{Rédaction, synthèse \& communication}

{\small
\cvtag{Rédaction technique}
\cvtag{Synthèse}
\cvtag{Communication orale}
\cvtag{Supports de présentation}
\cvtag{Comptes rendus}
\cvtag{Restitution}
\cvtag{Vulgarisation}
\cvtag{Recommandations}
\cvtag{Relations partenaires}
\cvtag{Anglais professionnel}
}

\divider\smallskip
\vspace{-.3cm}

\cvsubsection{R\&D, méthode \& systèmes complexes}

{\small
\cvtag{R\&D}
\cvtag{Validation}
\cvtag{Analyse de risques}
\cvtag{Simulation}
\cvtag{Modélisation}
\cvtag{Machine Learning}
\cvtag{Data Science}
\cvtag{Traitement du signal}
\cvtag{Systèmes complexes}
\cvtag{Traçabilité}
\cvtag{Documentation}
}

\divider\smallskip
\vspace{-.3cm}

\cvsubsection{Outils}

{\small
\cvtag{Python}
\cvtag{SQL}
\cvtag{Power BI}
\cvtag{Excel}
\cvtag{PowerPoint}
\cvtag{Git / GitLab}
\cvtag{Linux}
\cvtag{Confluence}
\cvtag{Jira}
\cvtag{Zenhub}
\cvtag{OpenSearch}
\cvtag{Sentry}
}

\end{document}
